% College Lecture 12 (asymmetric-key cryptography, RSA) + Forouzan 10.1--10.2. Every number recomputed by script
% (scratchpad/rsa/verify.py, gen.py): key tables, square-and-multiply chains, CRT, attack demos.
\section{Asymmetric-Key Cryptography and RSA \new}\label{sec:rsa}
\likely{New this year, so expect one Forouzan-style numerical: ``$p, q, e$ given: find $n$, $\phi(n)$, $d$; encrypt $P$; decrypt $C$'' (4--6 marks), plus a ``why'' question (proof of correctness, why $e = 2$ fails, one attack).}

\subsection{Why asymmetric-key cryptography exists \new}
\likely{``Differentiate symmetric-key and asymmetric-key cryptography'' / ``Why is symmetric cryptography still needed?'' (2--3 marks).}

In symmetric-key cryptography both parties use the same secret key, so before Alice and Bob can talk privately they must already share a secret. This is the \textbf{key distribution problem}: the key has to travel over a secure channel that usually does not exist, and in a group of $N$ users every pair needs its own key, so $N(N-1)/2$ secret keys must be created and stored (4950 for 100 users). Symmetric cryptography also cannot prove who sent a message, because anyone who can decrypt can also encrypt.

\textbf{Asymmetric-key} (public-key) cryptography removes the shared secret. Bob generates a \textbf{key pair}: a \textbf{public key} that he announces to everyone and a \textbf{private key} that only he knows. \imp{The public key locks; the private key unlocks.} Anyone can encrypt for Bob with his public key, but only his private key decrypts, so Alice needs no earlier contact with Bob and $N$ users need only $N$ key pairs. Formally $C = f(K_{\text{public}}, P)$ and $P = g(K_{\text{private}}, C)$, where $f$ is used only for encryption and $g$ only for decryption, and knowing $f$ must not reveal $g$. The secret is now carried by the receiver alone; for Bob to reply, he uses \emph{Alice's} public key. Plaintext and ciphertext are treated as \textbf{numbers}, and encryption is a number-theoretic function rather than substitution and permutation of symbols.

Asymmetric cryptography does \textbf{not} replace symmetric cryptography, because it is far slower: one RSA operation is a modular exponentiation with 1024--2048-bit numbers, while AES needs a few table lookups and XORs per byte (RSA is roughly a thousand times slower). It is therefore used where symmetric cryptography cannot work: \textbf{exchanging a secret session key, authentication and digital signatures}. In practice the two are combined: a random AES key is sent encrypted under the receiver's RSA key, and the bulk data is encrypted with AES under that key.

\begin{center}\small
\begin{tabular}{@{}lll@{}}
\toprule
 & Symmetric-key & Asymmetric-key\\
\midrule
Keys & one shared secret key & public/private key pair\\
Keys for $N$ users & $N(N-1)/2$ & $N$ pairs\\
Secret held by & both parties & receiver only\\
Based on & substitution, permutation & number theory (factoring, discrete log)\\
Speed & fast (bulk data) & slow (keys, signatures)\\
Examples & DES, AES & RSA, ElGamal, ECC\\
\bottomrule
\end{tabular}
\end{center}

\subsection{Trapdoor one-way functions \new}
\likely{``Define a trapdoor one-way function; give the RSA example (2 marks).}
The design problem is to publish $f$ while keeping $f^{-1}$ hidden. A \textbf{one-way function} (OWF) is a function $y = f(x)$ that is easy to compute, while computing $x = f^{-1}(y)$ from $y$ is computationally infeasible. A \textbf{trapdoor one-way function} (TOWF) has a third property: given $y$ \emph{and a secret} (the trapdoor), $x$ is easy to compute. A public-key cipher needs exactly this: the public key lets everyone evaluate $f$, and the trapdoor is the private key.

\textbf{Multiplication is one-way.} For large primes $p, q$, computing $n = p \times q$ is easy, but recovering $p$ and $q$ from $n$ is the \textbf{factorization problem}, for which no polynomial-time (classical) algorithm is known.

\textbf{$y = x^k \bmod n$ is a trapdoor one-way function.} Given $x$, $k$ and $n$, fast exponentiation (\S\ref{sec:modexp}) computes $y$ in about $2\log_2 k$ multiplications; given only $y$, $k$ and $n$, taking the $k$-th root modulo $n$ is believed to be as hard as factoring $n$. The trapdoor is $k'$ with $k k' \equiv 1 \pmod{\phi(n)}$: then $x = y^{k'} \bmod n$. Computing $k'$ needs $\phi(n) = (p-1)(q-1)$, and computing $\phi(n)$ needs $p$ and $q$, so the trapdoor exists only for whoever chose the primes. RSA is this function with $k = e$ and $k' = d$.

\subsection{RSA: key generation, encryption, decryption \new}
\likely{Write the three algorithms, then run them on small primes.}

RSA (Rivest, Shamir, Adleman, 1977) is the most widely used public-key cryptosystem. Bob builds the trapdoor from two primes and publishes only their product; Alice encrypts by raising the plaintext to the public exponent, and Bob decrypts by raising the ciphertext to the private exponent.

\begin{mustknow}[: RSA]
\textbf{Key generation (Bob)}
\begin{enumerate}[leftmargin=1.6em,itemsep=0pt,topsep=1pt]
  \item Choose two large primes $p \ne q$; compute $n = p \times q$ and $\phi(n) = (p-1)(q-1)$.
  \item Choose $e$ with $1 < e < \phi(n)$ and $\gcd(e, \phi(n)) = 1$.
  \item Compute $d = e^{-1} \bmod \phi(n)$ with the extended Euclidean algorithm, so $e \times d \equiv 1 \pmod{\phi(n)}$.
  \item Public key $(e, n)$, announced. Private key $d$, kept secret (with $p$, $q$, $\phi(n)$).
\end{enumerate}
\textbf{Encryption (Alice)}: $C = P^e \bmod n$, \;with $0 \le P < n$. \qquad \textbf{Decryption (Bob)}: $P = C^d \bmod n$.
\end{mustknow}

Encryption and decryption are each one fast exponentiation, so both are polynomial in the size of $n$. Eve sees $C$, $e$ and $n$ and must compute $P = \sqrt[e]{C} \bmod n$, which is exponential for her unless she can factor $n$. The plaintext must be smaller than $n$, because the result is only known modulo $n$: two plaintexts that differ by $n$ give the same ciphertext. A long message is cut into blocks, each a number less than $n$.

\subsection{Why decryption works: proof of correctness \new}
\likely{``Prove that RSA decryption recovers the plaintext'' (2--3 marks).}

Since $d$ is the inverse of $e$ modulo $\phi(n)$, $ed = 1 + k\,\phi(n)$ for some integer $k$. Decryption computes
\[
C^d \equiv (P^e)^d = P^{ed} = P^{1 + k\phi(n)} = P \cdot \big(P^{\phi(n)}\big)^k \pmod n.
\]
\textbf{Case 1: $\gcd(P, n) = 1$} (almost every $P$). Euler's theorem (\S\ref{sec:fermat}) gives $P^{\phi(n)} \equiv 1 \pmod n$, so $C^d \equiv P \cdot 1^k = P \pmod n$. Since $0 \le P < n$, the reduced result is exactly $P$, not just something congruent to it.

\textbf{Case 2: $\gcd(P, n) \ne 1$.} Euler's theorem does not apply, but the result still holds. Since $0 < P < pq$, $P$ is a multiple of exactly one prime, say $p \mid P$ and $q \nmid P$. Modulo $p$, $P \equiv 0$, so $P^{ed} \equiv 0 \equiv P$. Modulo $q$, Fermat's little theorem gives $P^{q-1} \equiv 1$, and $ed = 1 + k(p-1)(q-1)$, so $P^{ed} = P \cdot (P^{q-1})^{k(p-1)} \equiv P \pmod q$. Since $P^{ed} \equiv P$ modulo both $p$ and $q$, and $p, q$ are coprime, the CRT (\S\ref{sec:primes}) gives $P^{ed} \equiv P \pmod{pq}$. (This is Euler's ``second version'', $P^{k\phi(n)+1} \equiv P \pmod n$ for $n = pq$ and every $P$.) Check: $n = 221 = 13 \cdot 17$, $e = 5$, $d = 77$, $P = 65 = 5 \cdot 13$: $65^5 \bmod 221 = 182$ and $182^{77} \bmod 221 = 65$ \checkmark. Such a $P$ is no practical concern, because finding one would give $\gcd(P, n) = p$ and factor $n$.

\textbf{The two algebraic structures.} RSA works in two places (\S\ref{sec:algebra}). Encryption and decryption use the \textbf{public ring} $R = \langle \mathbb{Z}_n, +, \times \rangle$: $P$ and $C$ can be any element of $\mathbb{Z}_n$ (Case 2 shows non-invertible elements also work), and exponentiation is repeated multiplication, which anyone can do since $n$ is public. Key generation uses the \textbf{private group} $G = \langle \mathbb{Z}_{\phi(n)}^*, \times \rangle$: $e$ must belong to this group, i.e.\ $\gcd(e, \phi(n)) = 1$, so that its inverse $d$ exists in it. Only Bob can build $G$, because knowing $\phi(n)$ is as good as knowing the factors: $p + q = n - \phi(n) + 1$, and $p, q$ are the roots of $x^2 - (p+q)x + n = 0$. For $n = 187$, $\phi = 160$: $p + q = 28$, $x^2 - 28x + 187 = 0$ gives $x = (28 \pm 6)/2 = 17, 11$.

\subsection{The slide examples \new}
\likely{Verify one row by hand, e.g.\ ``with $n = 77$, $e = 13$, find $d$ and encrypt $5$.}
The lecture uses the key $p = 7$, $q = 11$, $e = 13$: $n = 77$, $\phi(n) = 60$, and $d = 37$ because $13 \times 37 = 481 = 8 \cdot 60 + 1$. Alice encrypts $5$ and John encrypts $63$ with the same public key $(13, 77)$; Bob decrypts both with $d = 37$. The table also has the class homework key ($n = 187$) and a small textbook key ($n = 33$); all values were recomputed.
\begin{center}\small
\begin{tabular}{@{}cccccll@{}}
\toprule
$p, q$ & $n$ & $\phi(n)$ & $e$ & $d$ & Encrypt & Decrypt\\
\midrule
7, 11 & 77 & 60 & 13 & 37 & $5^{13} \bmod 77 = 26$ & $26^{37} \bmod 77 = 5$\\
7, 11 & 77 & 60 & 13 & 37 & $63^{13} \bmod 77 = 28$ & $28^{37} \bmod 77 = 63$\\
17, 11 & 187 & 160 & 7 & 23 & $88^{7} \bmod 187 = 11$ & $11^{23} \bmod 187 = 88$\\
3, 11 & 33 & 20 & 7 & 3 & $5^7 \bmod 33 = 14$ & $14^3 \bmod 33 = 5$\\
\bottomrule
\end{tabular}
\end{center}
For the slide key, the $d$-table (\S\ref{sec:exteuclid}) on $(60, 13)$ has quotients $4, 1, 1, 1, 1, 2$ and ends with $r_1 = 1$, $t_1 = -23$, so $d = -23 + 60 = 37$. A negative $t_1$ is normal; add $\phi(n)$, never $n$.

\subsection{Exam-format RSA numericals \hot}
\asked{End-sem 2025--26 Q1(b) (6 marks): $M = 63$, $n = 77$, $e = 13$; find $d$ and the ciphertext by fast exponentiation ($d = 37$, $C = 28$).}
\likely{Exactly this layout: $\phi(n)$, gcd check, $d$-table, two square-and-multiply tables, a check line.}

Every exponentiation below uses repeated squaring (\S\ref{sec:modexp}): write the exponent as a sum of powers of 2, square repeatedly (reducing mod $n$ each time), then multiply the needed powers, reducing after each product.

\begin{worked}[: Stallings/Forouzan, the class homework key]
\textbf{Q.} In RSA, $p = 17$, $q = 11$, $e = 7$. Find $n$, $\phi(n)$ and the private key $d$. Encrypt $P = 88$ and show that decryption recovers it.

\textbf{Keys.} $n = 17 \times 11 = 187$; \;$\phi(n) = 16 \times 10 = 160$. Since $160 = 2^5 \cdot 5$ and $7$ is prime, $\gcd(7, 160) = 1$, so $e = 7$ is valid. $d = 7^{-1} \bmod 160$:
\begin{center}\small
\begin{tabular}{@{}r|rrr|rrr@{}}
\toprule
$q$ & $r_1$ & $r_2$ & $r$ & $t_1$ & $t_2$ & $t$\\
\midrule
22 & 160 & 7 & 6 & 0 & 1 & $-22$\\
1 & 7 & 6 & 1 & 1 & $-22$ & 23\\
6 & 6 & 1 & 0 & $-22$ & 23 & $-160$\\
 & \textbf{1} & 0 & & \textbf{23} & $-160$ & \\
\bottomrule
\end{tabular}
\end{center}
$r_1 = 1$ (inverse exists), $d = t_1 = \mathbf{23}$. Check: $7 \times 23 = 161 = 160 + 1$ \checkmark. Public key $(7, 187)$, private key $23$.

\textbf{Encrypt} $C = 88^7 \bmod 187$, with $7 = 4 + 2 + 1$:
\begin{center}\small
\begin{tabular}{@{}lll@{}}
\toprule
Power & Computation & mod 187\\
\midrule
$88^{1}$ & & 88\\
$88^{2}$ & $88^2 = 7744 = 41\cdot187 + 77$ & 77\\
$88^{4}$ & $77^2 = 5929 = 31\cdot187 + 132$ & 132\\
\bottomrule
\end{tabular}
\end{center}
$88^7 = 88^4 \cdot 88^2 \cdot 88^1$: $132 \cdot 77 = 10164 \equiv 66$; \;$66 \cdot 88 = 5808 \equiv 11$. \;$\Rightarrow C = \mathbf{11}$.

\textbf{Decrypt} $P = 11^{23} \bmod 187$, with $23 = 16 + 4 + 2 + 1$:
\begin{center}\small
\begin{tabular}{@{}lll@{}}
\toprule
Power & Computation & mod 187\\
\midrule
$11^{1}$ & & 11\\
$11^{2}$ & $11^2 = 121$ & 121\\
$11^{4}$ & $121^2 = 14641 = 78\cdot187 + 55$ & 55\\
$11^{8}$ & $55^2 = 3025 = 16\cdot187 + 33$ & 33\\
$11^{16}$ & $33^2 = 1089 = 5\cdot187 + 154$ & 154\\
\bottomrule
\end{tabular}
\end{center}
$154 \cdot 55 = 8470 \equiv 55$; \;$55 \cdot 121 = 6655 \equiv 110$; \;$110 \cdot 11 = 1210 \equiv 88$. \;$\Rightarrow P = \mathbf{88}$ \checkmark\ (the original plaintext).
\end{worked}

\begin{worked}[: new, Forouzan exercise style (only $n$ given)]
\textbf{Q.} Bob's RSA public key is $(e, n) = (5, 221)$. (a) Find his private key. (b) Alice sends $P = 100$; find $C$. (c) Show Bob's decryption.

\textbf{(a)} Eve cannot factor a 1024-bit $n$, but $221$ factors by trial division: $221 = 13 \times 17$. So $\phi(n) = 12 \times 16 = 192 = 2^6 \cdot 3$, and $\gcd(5, 192) = 1$. $d = 5^{-1} \bmod 192$:
\begin{center}\small
\begin{tabular}{@{}r|rrr|rrr@{}}
\toprule
$q$ & $r_1$ & $r_2$ & $r$ & $t_1$ & $t_2$ & $t$\\
\midrule
38 & 192 & 5 & 2 & 0 & 1 & $-38$\\
2 & 5 & 2 & 1 & 1 & $-38$ & 77\\
2 & 2 & 1 & 0 & $-38$ & 77 & $-192$\\
 & \textbf{1} & 0 & & \textbf{77} & $-192$ & \\
\bottomrule
\end{tabular}
\end{center}
$d = \mathbf{77}$. Check: $5 \times 77 = 385 = 2 \cdot 192 + 1$ \checkmark.

\textbf{(b)} $5 = 4 + 1$: \;$100^2 = 10000 = 45 \cdot 221 + 55$; \;$55^2 = 3025 = 13 \cdot 221 + 152$, so $100^4 \equiv 152$. \;$C = 152 \cdot 100 = 15200 = 68 \cdot 221 + 172$ $\Rightarrow C = \mathbf{172}$.

\textbf{(c)} $P = 172^{77} \bmod 221$, with $77 = 64 + 8 + 4 + 1$:
\begin{center}\small
\begin{tabular}{@{}lll@{}}
\toprule
Power & Computation & mod 221\\
\midrule
$172^{1}$ & & 172\\
$172^{2}$ & $172^2 = 29584 = 133\cdot221 + 191$ & 191\\
$172^{4}$ & $191^2 = 36481 = 165\cdot221 + 16$ & 16\\
$172^{8}$ & $16^2 = 256 = 221 + 35$ & 35\\
$172^{16}$ & $35^2 = 1225 = 5\cdot221 + 120$ & 120\\
$172^{32}$ & $120^2 = 14400 = 65\cdot221 + 35$ & 35\\
$172^{64}$ & $35^2 \equiv 120$ (as above) & 120\\
\bottomrule
\end{tabular}
\end{center}
$172^{64} \cdot 172^{8} = 120 \cdot 35 = 4200 = 19 \cdot 221 + 1 \equiv 1$; \;$1 \cdot 16 = 16$; \;$16 \cdot 172 = 2752 = 12 \cdot 221 + 100$. \;$\Rightarrow P = \mathbf{100}$ \checkmark.
\end{worked}

\subsection{Choosing $e$; why $e = 2$ never works \new}
\likely{``Why is $e = 2$ not allowed? or ``is $e = 3$ valid for $p = 19$, $q = 23$? (1--2 marks).}
Encryption cost depends on the number of bits and of 1-bits in $e$, so a small $e$ with few 1-bits makes encryption cheap. The standard choice is $e = 2^{16} + 1 = 65537$: it is prime (so it is coprime to $\phi(n)$ unless a prime factor of $\phi(n)$ is 65537 itself), its binary form $1\,0000\,0000\,0000\,0001$ needs only 16 squarings and one multiplication, and it is large enough to defeat the low-exponent attacks below. The private exponent $d$ is then about as long as $n$, so decryption is the slow side (which the CRT speeds up).

\textbf{$e = 2$ (or any even $e$) is invalid.} Since $p$ and $q$ are odd primes, $p - 1$ and $q - 1$ are even, so $\phi(n)$ is even (in fact divisible by 4). Then $\gcd(2, \phi(n)) = 2 \ne 1$, $e$ is not in $\mathbb{Z}_{\phi(n)}^*$, and no $d$ exists. The cause is that squaring is not one-to-one: $x$ and $n - x$ have the same square, so $C = P^2$ cannot be uniquely undone. The same check catches a Forouzan exercise: $p = 19$, $q = 23$, $e = 3$ gives $\phi(n) = 18 \cdot 22 = 396 = 3 \cdot 132$, so $\gcd(3, 396) = 3$ and $e = 3$ is not allowed for this $n$.

\begin{trap}[: the usual RSA slips]
$d$ is the inverse of $e$ modulo $\boldsymbol{\phi(n)}$, never modulo $n$ (for $(7, 187)$ the answer is $23$; $7^{-1} \bmod 187 = 107$ is useless). Check $\gcd(e, \phi(n)) = 1$ before running the table. The plaintext must satisfy $P < n$. Encrypt with the \emph{receiver's} public key. Reduce mod $n$ after every multiplication, and turn a negative $t_1$ into $t_1 + \phi(n)$.
\end{trap}

\subsection{Faster decryption with the CRT \new}
\likely{``Decrypt using the Chinese remainder theorem'' on a small key, or ``why is CRT decryption about 4 times faster''.}

Decryption is slow because $d$ and $n$ are both full size. Bob, however, knows $p$ and $q$, so he can work modulo each prime separately and combine the answers with the CRT (\S\ref{sec:primes}). Modulo $p$ the exponent can be reduced: since $C^{p-1} \equiv 1 \pmod p$ (Fermat), $C^d \equiv C^{d \bmod (p-1)} \pmod p$. Bob computes
\[
d_p = d \bmod (p-1),\quad d_q = d \bmod (q-1),\qquad P_p = C^{d_p} \bmod p,\quad P_q = C^{d_q} \bmod q,
\]
and then the unique $P < n$ with $P \equiv P_p \pmod p$ and $P \equiv P_q \pmod q$:
\[
P = \big(P_p \cdot q \cdot (q^{-1} \bmod p) + P_q \cdot p \cdot (p^{-1} \bmod q)\big) \bmod n .
\]
Each half uses a modulus of half the length (a multiplication costs about $1/4$) and an exponent of half the length (half as many steps), so each half costs about $1/8$ of the full exponentiation and both together about $1/4$: decryption is roughly \textbf{4 times faster}. Only the key owner can use it, since it needs $p$ and $q$.

\begin{worked}[: CRT decryption of the previous ciphertext]
\textbf{Q.} With $p = 13$, $q = 17$, $d = 77$, decrypt $C = 172$ using the CRT.

$d_p = 77 \bmod 12 = 5$, \;$d_q = 77 \bmod 16 = 13$; \;$C_p = 172 \bmod 13 = 3$, \;$C_q = 172 \bmod 17 = 2$.\\
$P_p = 3^5 \bmod 13 = 243 \bmod 13 = 9$. \;$P_q = 2^{13} \bmod 17$: $2^4 = 16 \equiv -1$, $2^8 \equiv 1$, so $2^{13} = 2^8 \cdot 2^4 \cdot 2 \equiv -2 \equiv 15$.\\
$q^{-1} \bmod p = 17^{-1} \bmod 13 = 4^{-1} \bmod 13 = 10$ \;($4 \cdot 10 = 40 = 3 \cdot 13 + 1$); \;$p^{-1} \bmod q = 13^{-1} \bmod 17 = 4$ \;($13 \cdot 4 = 52 = 3 \cdot 17 + 1$).\\
$P = 9 \cdot 17 \cdot 10 + 15 \cdot 13 \cdot 4 = 1530 + 780 = 2310 = 10 \cdot 221 + 100$ $\Rightarrow P = \mathbf{100}$ \checkmark, the same as $172^{77} \bmod 221$, with exponents 5 and 13 instead of 77.
\end{worked}

\subsection{Attacks on RSA \new}
\likely{``Explain any two attacks on RSA'' or ``why must $p$ and $q$ be large / not close / why OAEP''.}

No devastating attack on correctly used RSA is known; every known attack exploits a weak choice of parameters or a careless use of textbook RSA. The slide's list, each with the reason it works:
\begin{itemize}[leftmargin=1.4em]
  \item \textbf{Factorization.} If Eve factors $n$ she computes $\phi(n)$ and then $d = e^{-1} \bmod \phi(n)$, exactly as Bob did. The whole security rests on $n$ being too large to factor.
  \item \textbf{Chosen-ciphertext.} RSA is multiplicative: $(P_1 P_2)^e = P_1^e P_2^e$. To decrypt $C$, Eve picks $X$, asks Bob to decrypt the innocent-looking $Y = C \cdot X^e \bmod n$, receives $Z = Y^d = P \cdot X$, and computes $P = Z \cdot X^{-1}$. With $(7, 187)$, $C = 11$, $X = 2$: $Y = 11 \cdot 128 \bmod 187 = 99$, $Z = 99^{23} \bmod 187 = 176$, $P = 176 \cdot 94 \bmod 187 = 88$ \checkmark.
  \item \textbf{Low encryption exponent.} If $e$ is small and $P^e < n$, no reduction happens and $P = \sqrt[e]{C}$ over the integers ($e = 3$, $n = 187$, $P = 5$: $C = 125$, $\sqrt[3]{125} = 5$). \textbf{Coppersmith's} theorem extends this to finding small unknown parts of $P$. \textbf{Broadcast}: if the same $P$ goes to $e$ receivers with moduli $n_1, \dots, n_e$, the CRT gives $P^e \bmod n_1 \cdots n_e$, which equals $P^e$ because $P^e < n_1 \cdots n_e$, so an integer root reveals $P$ (with $e = 3$, $n = 55, 391, 1189$ and $P = 42$: ciphertexts $3, 189, 370$, CRT gives $74088 = 42^3$). \textbf{Related messages} (Franklin--Reiter) and \textbf{short pad} attacks similarly recover messages that are linearly related or padded with too few random bits.
  \item \textbf{Low decryption exponent.} A small $d$ speeds decryption but Wiener's attack recovers $d$ from $e/n$ by continued fractions when $d < \tfrac{1}{3} n^{1/4}$. A \textbf{revealed} $d$ is worse than it looks: $n$ with one pair $(e, d)$ lets anyone factor $n$, so all keys on that modulus are dead.
  \item \textbf{Plaintext attacks.} \emph{Short message}: if $P$ comes from a small set (a PIN), Eve encrypts every candidate with the public key and compares, since textbook RSA is deterministic. \emph{Cyclic}: $C$ is a permutation of $\mathbb{Z}_n$, so re-encrypting $C, C^e, C^{e^2}, \dots$ eventually returns to $C$, and the value just before $C$ is $P$; with $(7, 33)$ and $C = 14$: $14 \to 20 \to 26 \to 5 \to 14$, so $P = 5$. \emph{Unconcealed}: some $P$ satisfy $P^e \equiv P$ (always $0$, $1$, $n - 1$); their number is $(1 + \gcd(e-1, p-1))(1 + \gcd(e-1, q-1))$, which for the slide key $(13, 77)$ is $7 \cdot 3 = 21$ messages (e.g.\ $10^{13} \equiv 10$), because $e - 1 = 12$ is a multiple of $p - 1 = 6$.
  \item \textbf{Common modulus.} If two users share $n$ with exponents $e_1, e_2$, $\gcd(e_1, e_2) = 1$, and the same $P$ is sent to both, Eve finds $s e_1 + t e_2 = 1$ and gets $C_1^s C_2^t = P^{s e_1 + t e_2} = P$. With $n = 187$, $e_1 = 3$, $e_2 = 7$, $P = 5$: $C_1 = 125$, $C_2 = 146$, $s = -2$, $t = 1$, $P = (125^{-1})^2 \cdot 146 = 3^2 \cdot 146 = 1314 \equiv 5$ \checkmark. Also, any user of the shared $n$ can factor it from his own $(e, d)$.
  \item \textbf{Implementation (timing, power).} Square-and-multiply does an extra multiplication for every 1-bit of $d$, so the decryption time or power trace leaks the bits of $d$. Defences: constant-time exponentiation, random delays, or \emph{blinding} (decrypt $C \cdot r^e$, then multiply by $r^{-1}$).
\end{itemize}

\subsection{Recommendations}
Each recommendation blocks one of the attacks above.
\begin{itemize}[leftmargin=1.4em]
  \item $n$ at least \textbf{1024 bits} (about 309 decimal digits; 2048 today), so each of $p$, $q$ at least \textbf{512 bits} (about 154 digits). This defeats factorization.
  \item $p$ and $q$ must not be close to each other, otherwise Fermat factoring ($n = a^2 - b^2$ with $a \approx \sqrt n$) finds them quickly. Both $p - 1$ and $q - 1$ should have a large prime factor (against Pollard's $p - 1$ method), and $p/q$ should not be close to a fraction with small numerator and denominator.
  \item The modulus must not be shared between users (common-modulus attack).
  \item $e = 2^{16} + 1$ or an integer close to it (against low-exponent attacks, while keeping encryption fast).
  \item If $d$ leaks, Bob must change $n$, $e$ and $d$, since one $(e, d)$ pair reveals the factors of $n$.
  \item Messages must be padded with \textbf{OAEP} (Optimal Asymmetric Encryption Padding). OAEP mixes the message with a fresh random value $r$ through two hash-based masks ($M_1 = P \oplus G(r)$, $M_2 = r \oplus H(M_1)$) and encrypts $M_1 \| M_2$. This makes RSA probabilistic (same $P$, different $C$ each time), fills the block so $P^e > n$, and destroys the multiplicative structure, blocking the short-message, broadcast, cyclic-style and chosen-ciphertext tricks.
\end{itemize}

\begin{fixbox}
\begin{itemize}[leftmargin=1.2em,itemsep=0pt]
  \item Slide on trapdoor functions: ``there is not a polynomial time solution'' to factoring should read \emph{no polynomial-time classical algorithm is known}; it is not proven impossible (Shor's quantum algorithm is polynomial).
  \item Slide on RSA size: a 512-bit prime has 154 or 155 decimal digits ($2^{511}$ has 154, $2^{512}$ has 155); 1024 bits is 309 digits, as stated.
\end{itemize}
\end{fixbox}
